
\documentclass[11pt]{article}

\usepackage[T1]{fontenc}
\usepackage[utf8]{inputenc}
\usepackage{lmodern}
\usepackage{amsmath,amssymb,amsfonts,bm,mathtools}
\usepackage[a4paper,margin=2.5cm]{geometry}
\usepackage{hyperref}
\newcommand{\kk}{k}

\newcommand{\bra}[1] { \langle #1 | }
\newcommand{\ket}[1] { | #1 \rangle }
\newcommand{\braket}[2] { \langle #1 | #2 \rangle }

\title{Detailed Derivation of the Hubbard Contribution to $\Delta \mathbf{M}$\\
and Its Analogy with $\Delta \mathbf{M}_{\mathrm{bare}}$}
\author{}
\date{}

\begin{document}

\maketitle

\section{Goal}

The purpose of this note is to derive step by step the contribution to the orbital magnetization associated with a Hubbard potential, starting from the same GIPAW-like expression used for the bare non-local contribution. The goal is to show that the final Hubbard term has exactly the same structure as $\Delta \mathbf{M}_{\mathrm{bare}}$, upon replacing non-local pseudopotential projectors with Hubbard projectors.

Throughout the derivation we assume the norm-conserving case, so that the overlap operator is the identity:
\begin{equation}
S = \mathbb{I}.
\end{equation}
Therefore no ultrasoft overlap corrections appear.

\section{Starting point}

We start from the standard bare contribution written in operator form:
\begin{equation}
\Delta \mathbf{M}
=
\frac{\alpha}{2}
\sum_I
\left\langle
(\mathbf{R}_I-\mathbf{r}) \times \frac{1}{i}
\left[
\mathbf{R}_I-\mathbf{r},\hat V_I
\right]
\right\rangle .
\label{eq:start}
\end{equation}

Here:
\begin{itemize}
\item $I$ labels the atomic site,
\item $\mathbf{R}_I$ is the position of site $I$,
\item $\hat V_I$ is an operator localized on site $I$.
\end{itemize}

For convenience, define
\begin{equation}
\bm{\rho}_I = \mathbf{r}-\mathbf{R}_I.
\label{eq:rho_def}
\end{equation}
Then
\begin{equation}
\mathbf{R}_I-\mathbf{r} = -\bm{\rho}_I,
\end{equation}
and
\begin{equation}
\left[
\mathbf{R}_I-\mathbf{r},\hat V_I
\right]
=
\left[
-\bm{\rho}_I,\hat V_I
\right]
=
-\left[
\bm{\rho}_I,\hat V_I
\right].
\end{equation}
Therefore Eq.~\eqref{eq:start} becomes
\begin{equation}
\Delta \mathbf{M}
=
\frac{\alpha}{2}
\sum_I
\left\langle
\bm{\rho}_I \times \frac{1}{i}
\left[
\bm{\rho}_I,\hat V_I
\right]
\right\rangle .
\label{eq:start_rho}
\end{equation}

\section{Expansion of the commutator}

We expand the commutator:
\begin{equation}
\left[
\bm{\rho}_I,\hat V_I
\right]
=
\bm{\rho}_I \hat V_I - \hat V_I \bm{\rho}_I .
\end{equation}
Substituting into Eq.~\eqref{eq:start_rho}, we obtain
\begin{equation}
\Delta \mathbf{M}
=
\frac{\alpha}{2}
\sum_I
\left\langle
\bm{\rho}_I \times \frac{1}{i}
\left(
\bm{\rho}_I \hat V_I - \hat V_I \bm{\rho}_I
\right)
\right\rangle .
\label{eq:expand_comm}
\end{equation}

Using linearity of the cross product,
\begin{equation}
\bm{\rho}_I \times
\left(
\bm{\rho}_I \hat V_I - \hat V_I \bm{\rho}_I
\right)
=
\bm{\rho}_I \times (\bm{\rho}_I \hat V_I)
-
\bm{\rho}_I \times (\hat V_I \bm{\rho}_I) .
\label{eq:cross_split}
\end{equation}

Let us analyze the first term. In components,
\begin{equation}
\left[\bm{\rho}_I \times (\bm{\rho}_I \hat V_I)\right]_\gamma
=
\sum_{\mu\nu}
\epsilon_{\gamma\mu\nu}
\,\rho_{I,\mu}\rho_{I,\nu}\hat V_I .
\end{equation}
Since $\rho_{I,\mu}\rho_{I,\nu}$ is symmetric under $\mu \leftrightarrow \nu$, while $\epsilon_{\gamma\mu\nu}$ is antisymmetric, this term vanishes:
\begin{equation}
\bm{\rho}_I \times (\bm{\rho}_I \hat V_I)=0 .
\end{equation}
Therefore,
\begin{equation}
\bm{\rho}_I \times
\left(
\bm{\rho}_I \hat V_I - \hat V_I \bm{\rho}_I
\right)
=
-
\bm{\rho}_I \times (\hat V_I \bm{\rho}_I) .
\end{equation}

Hence Eq.~\eqref{eq:expand_comm} becomes
\begin{equation}
\Delta \mathbf{M}
=
-\frac{\alpha}{2}
\sum_I
\left\langle
\bm{\rho}_I \times \frac{1}{i}\hat V_I \bm{\rho}_I
\right\rangle .
\label{eq:generic_form}
\end{equation}

Equivalently, since $1/i=-i$,
\begin{equation}
\Delta \mathbf{M}
=
\frac{i\alpha}{2}
\sum_I
\left\langle
\bm{\rho}_I \times \hat V_I \bm{\rho}_I
\right\rangle .
\label{eq:generic_form_i}
\end{equation}


\section{Recovery of the bare non-local formula}

Consider first a separable non-local pseudopotential:
\begin{equation}
\hat V_I^{\mathrm{NL}}
=
\sum_{ij}
|\beta_i^I\rangle D_{ij}^I \langle \beta_j^I| .
\label{eq:vnl_sep}
\end{equation}

Substituting Eq.~\eqref{eq:vnl_sep} into Eq.~\eqref{eq:generic_form}, we obtain
\begin{equation}
\Delta \mathbf{M}_{\mathrm{bare}}
=
-\frac{\alpha}{2}
\sum_I
\left\langle
\bm{\rho}_I \times \frac{1}{i}
\sum_{ij}
|\beta_i^I\rangle D_{ij}^I \langle \beta_j^I|
\bm{\rho}_I
\right\rangle .
\end{equation}
By linearity,
\begin{equation}
\Delta \mathbf{M}_{\mathrm{bare}}
=
-\frac{\alpha}{2}
\sum_I \sum_{ij}
\left\langle
\bm{\rho}_I \times \frac{1}{i}
|\beta_i^I\rangle D_{ij}^I \langle \beta_j^I|
\bm{\rho}_I
\right\rangle .
\end{equation}

Evaluating the expectation value over occupied Bloch states gives
\begin{equation}
\Delta \mathbf{M}_{\mathrm{bare}}
=
-\frac{\alpha}{2}
\sum_{n\mathbf{k}}^{\mathrm{occ}} f_{n\mathbf{k}}
\sum_I \sum_{ij}
\left\langle u_{n\mathbf{k}} \left|
\bm{\rho}_I \times \frac{1}{i}
|\beta_i^{\mathbf{k},I}\rangle D_{ij}^I
\langle \beta_j^{\mathbf{k},I}|
\bm{\rho}_I
\right| u_{n\mathbf{k}} \right\rangle .
\label{eq:bare_exp}
\end{equation}

Using
\begin{equation}
(\mathbf{A}\times\mathbf{B})_\gamma
=
\sum_{\mu\nu}
\epsilon_{\gamma\mu\nu} A_\mu B_\nu ,
\label{eq:cross_components}
\end{equation}
we obtain
\begin{equation}
\Delta M_{\mathrm{bare},\gamma}
=
-\frac{\alpha}{2}
\sum_{n\mathbf{k}}^{\mathrm{occ}} f_{n\mathbf{k}}
\sum_I \sum_{ij} \sum_{\mu\nu}
\epsilon_{\gamma\mu\nu}
\frac{1}{i}
\langle u_{n\mathbf{k}} | \rho_{I,\mu} | \beta_i^{\mathbf{k},I} \rangle
D_{ij}^I
\langle \beta_j^{\mathbf{k},I} | \rho_{I,\nu} | u_{n\mathbf{k}} \rangle .
\label{eq:bare_components}
\end{equation}

More precisely, for a projector localized on atom $I$, we use the
no-phase KB projector, i.e. the projector in which the global
$e^{-i\mathbf{k}\cdot\mathbf{R}_I}$ phase associated with the atomic
center has been removed. This convention is discussed in
Appendix~\ref{app:kb_no_phase}. With this choice, the dependence on the
local coordinate
\begin{equation}
\boldsymbol{\rho}_I = \mathbf r-\mathbf R_I
\end{equation}
is made explicit by writing
\begin{equation}
\widetilde{\beta}_i^{\mathbf{k},I}(\mathbf{r})
=
e^{-i\mathbf{k}\cdot(\mathbf{r}-\mathbf{R}_I)}
\beta_i^{I}(\mathbf{r}),
\end{equation}
where $\beta_i^{I}(\mathbf r)$ denotes the localized projector already
centered on atom $I$.

It follows that
\begin{equation}
\partial_{k_\mu}\widetilde{\beta}_i^{\mathbf{k},I}(\mathbf{r})
=
-i (r_\mu-R_{I,\mu})
\widetilde{\beta}_i^{\mathbf{k},I}(\mathbf{r})
=
-i \rho_{I,\mu}
\widetilde{\beta}_i^{\mathbf{k},I}(\mathbf{r}).
\end{equation}
Hence,
\begin{equation}
\rho_{I,\mu}\widetilde{\beta}_i^{\mathbf{k},I}
=
i\,\partial_{k_\mu}\widetilde{\beta}_i^{\mathbf{k},I}.
\label{eq:rho_beta_ket}
\end{equation}
This gives
\begin{equation}
\langle u_{n\mathbf{k}} | \rho_{I,\mu} |
\widetilde{\beta}_i^{\mathbf{k},I} \rangle
=
i \,
\langle u_{n\mathbf{k}} |
\partial_{k_\mu} \widetilde{\beta}_i^{\mathbf{k},I} \rangle ,
\label{eq:rho_beta_left}
\end{equation}
and, by Hermitian conjugation,
\begin{equation}
\langle \widetilde{\beta}_j^{\mathbf{k},I} |
\rho_{I,\nu} | u_{n\mathbf{k}} \rangle
=
- i \,
\langle
\partial_{k_\nu} \widetilde{\beta}_j^{\mathbf{k},I}
| u_{n\mathbf{k}} \rangle .
\label{eq:rho_beta_right}
\end{equation}



Substituting into Eq.~\eqref{eq:bare_components},
\begin{align}
\Delta M_{\mathrm{bare},\gamma}
&=
-\frac{\alpha}{2}
\sum_{n\mathbf{k}}^{\mathrm{occ}} f_{n\mathbf{k}}
\sum_I \sum_{ij} \sum_{\mu\nu}
\epsilon_{\gamma\mu\nu}
\frac{1}{i}
\left(
i \langle u_{n\mathbf{k}} | \partial_{k_\mu}
\widetilde{\beta}_i^{\mathbf{k},I} \rangle
\right)
D_{ij}^I
\left(
- i \langle \partial_{k_\nu}
\widetilde{\beta}_j^{\mathbf{k},I} | u_{n\mathbf{k}} \rangle
\right)
\nonumber\\
&=
\frac{i\alpha}{2}
\sum_{n\mathbf{k}}^{\mathrm{occ}} f_{n\mathbf{k}}
\sum_I \sum_{ij} \sum_{\mu\nu}
\epsilon_{\gamma\mu\nu}
\langle u_{n\mathbf{k}} | \partial_{k_\mu}
\widetilde{\beta}_i^{\mathbf{k},I} \rangle
D_{ij}^I
\langle \partial_{k_\nu}
\widetilde{\beta}_j^{\mathbf{k},I} | u_{n\mathbf{k}} \rangle .
\label{eq:bare_before_im}
\end{align}

Since the contraction with $\epsilon_{\gamma\mu\nu}$ selects the
antisymmetric part in the Cartesian indices, the result can be written
in manifestly real form as
\begin{equation}
\Delta M_{\mathrm{bare},\gamma}
=
-\frac{\alpha}{2}
\sum_{n\mathbf{k}}^{\mathrm{occ}} f_{n\mathbf{k}}
\sum_I \sum_{ij} \sum_{\mu\nu}
\epsilon_{\gamma\mu\nu}
\,\mathrm{Im}
\left[
\langle u_{n\mathbf{k}} | \partial_{k_\mu}
\widetilde{\beta}_i^{\mathbf{k},I} \rangle
D_{ij}^I
\langle \partial_{k_\nu}
\widetilde{\beta}_j^{\mathbf{k},I} | u_{n\mathbf{k}} \rangle
\right] .
\label{eq:deltam_bare_final}
\end{equation}

\noindent\textit{Note.}
In the implementation the full double sum over $(\mu,\nu)$ is not
performed explicitly. For a given magnetization component $\gamma$, only
one cyclic pair $(\mu,\nu)$ is evaluated. For instance,
\[
(\gamma,\mu,\nu)=(1,2,3),(2,3,1),(3,1,2),
\]
so that $\epsilon_{\gamma\mu\nu}=+1$.

Let
\begin{equation}
Z_{\mu\nu}
=
\frac{\alpha}{2}
\sum_{n\mathbf{k}}^{\mathrm{occ}} f_{n\mathbf{k}}
\sum_I \sum_{ij}
\langle u_{n\mathbf{k}} | \partial_{k_\mu}
\widetilde{\beta}_i^{\mathbf{k},I} \rangle
D_{ij}^I
\langle \partial_{k_\nu}
\widetilde{\beta}_j^{\mathbf{k},I} | u_{n\mathbf{k}} \rangle .
\end{equation}
Using the Hermiticity of $D_{ij}^I$, one has
\begin{equation}
Z_{\nu\mu}=Z_{\mu\nu}^{*}.
\end{equation}
Therefore the two antisymmetric terms selected by
$\epsilon_{\gamma\mu\nu}$ give
\begin{align}
\epsilon_{\gamma\mu\nu}\mathrm{Im}\,Z_{\mu\nu}
+
\epsilon_{\gamma\nu\mu}\mathrm{Im}\,Z_{\nu\mu}
&=
\mathrm{Im}\,Z_{\mu\nu}
-
\mathrm{Im}\,Z_{\nu\mu}
\nonumber\\
&=
\mathrm{Im}\,Z_{\mu\nu}
-
\mathrm{Im}\,Z_{\mu\nu}^{*}
\nonumber\\
&=
2\,\mathrm{Im}\,Z_{\mu\nu}.
\end{align}
Thus, evaluating only one cyclic Cartesian pair requires multiplying the
result by a factor of two. With the overall negative sign in
Eq.~\eqref{eq:deltam_bare_final}, the implemented contribution is
therefore proportional to
\begin{equation}
-2\,\mathrm{Im}\,Z_{\mu\nu}.
\end{equation}
This corresponds to the code statement
\begin{verbatim}
delta_M_bare(kk) = delta_M_bare(kk) - 2.d0*imag(tmp)
\end{verbatim}
where \texttt{tmp} is the numerical realization of $Z_{\mu\nu}$ for the
cyclic pair associated with the component \texttt{kk}.

\section{Hubbard potential in separable form}

In the norm-conserving case, the Hubbard potential localized on site $I$
can be written in separable form as
\begin{equation}
\hat V_I^{\mathrm{Hub}}
=
\sum_{mm'}
|\phi_{Im}^{\mathbf{k}}\rangle
V_{mm'}^I
\langle \phi_{Im'}^{\mathbf{k}}| .
\label{eq:vhub_sep}
\end{equation}
Here:
\begin{itemize}
\item $|\phi_{Im}^{\mathbf{k}}\rangle$ and
$\langle \phi_{Im'}^{\mathbf{k}}|$ denote the Bloch-summed Hubbard
projectors associated with site $I$;
\item depending on the chosen Hubbard-projector scheme, these projectors
may be atomic orbitals, orthogonalized atomic orbitals, or projectors
read from the pseudopotential file;
\item $m,m'$ label the local Hubbard manifold;
\item $V_{mm'}^I$ is the Hubbard potential matrix on site $I$.
\end{itemize}

\section{Derivation of the Hubbard contribution}
\label{sec:hubmagnetization}

We now repeat the same derivation with $\hat V_I = \hat V_I^{\mathrm{Hub}}$. Starting from Eq.~\eqref{eq:generic_form},
\begin{equation}
\Delta \mathbf{M}_{\mathrm{Hub}}
=
-\frac{\alpha}{2}
\sum_I
\left\langle
\bm{\rho}_I \times \frac{1}{i}\hat V_I^{\mathrm{Hub}} \bm{\rho}_I
\right\rangle .
\label{eq:hub_start}
\end{equation}

Substituting Eq.~\eqref{eq:vhub_sep},
\begin{equation}
\Delta \mathbf{M}_{\mathrm{Hub}}
=
-\frac{\alpha}{2}
\sum_I
\left\langle
\bm{\rho}_I \times \frac{1}{i}
\sum_{mm'}
V_{mm'}^I
{\phi_{Im}^{\mathbf{k}}}{\phi_{Im'}^{\mathbf{k}}}
\bm{\rho}_I
\right\rangle .
\end{equation}
By linearity,
\begin{equation}
\Delta \mathbf{M}_{\mathrm{Hub}}
=
-\frac{\alpha}{2}
\sum_I \sum_{mm'}
\left\langle
\bm{\rho}_I \times \frac{1}{i}
V_{mm'}^I
{\phi_{Im}^{\mathbf{k}}}{\phi_{Im'}^{\mathbf{k}}}
\bm{\rho}_I
\right\rangle .
\end{equation}

Evaluating over occupied Bloch states,
\begin{equation}
\Delta \mathbf{M}_{\mathrm{Hub}}
=
-\frac{\alpha}{2}
\sum_{n\mathbf{k}}^{\mathrm{occ}} f_{n\mathbf{k}}
\sum_I \sum_{mm'}
\left\langle u_{n\mathbf{k}} \left|
\bm{\rho}_I \times \frac{1}{i}
V_{mm'}^I
{\phi_{Im}^{\mathbf{k}}}{\phi_{Im'}^{\mathbf{k}}}
\bm{\rho}_I
\right| u_{n\mathbf{k}} \right\rangle .
\label{eq:hub_exp}
\end{equation}

Using Eq.~\eqref{eq:cross_components},
\begin{equation}
\Delta M_{\mathrm{Hub},\gamma}
=
-\frac{\alpha}{2}
\sum_{n\mathbf{k}}^{\mathrm{occ}} f_{n\mathbf{k}}
\sum_I \sum_{mm'} \sum_{\mu\nu}
\epsilon_{\gamma\mu\nu}
\frac{1}{i}
\langle u_{n\mathbf{k}} | \rho_{I,\mu} | \phi_{Im}^{\mathbf{k}} \rangle
V_{mm'}^I
\langle \phi_{Im'}^{\mathbf{k}} | \rho_{I,\nu} | u_{n\mathbf{k}} \rangle .
\label{eq:hub_components}
\end{equation}

In complete analogy with the KB projectors, the Hubbard projectors
entering the $k$ derivatives are understood to be no-phase projectors.
We denote them by $\widetilde{\phi}_{Im}^{\mathbf{k}}$. The tilde will
be omitted in the final compact expressions only when no ambiguity is
possible.

Now use the analogues of Eqs.~\eqref{eq:rho_beta_left} and
\eqref{eq:rho_beta_right} for the no-phase Hubbard projectors:
\begin{equation}
\langle u_{n\mathbf{k}} | \rho_{I,\mu} |
\widetilde{\phi}_{Im}^{\mathbf{k}} \rangle
=
i \,
\langle u_{n\mathbf{k}} |
\partial_{k_\mu} \widetilde{\phi}_{Im}^{\mathbf{k}} \rangle ,
\label{eq:rho_phi_left}
\end{equation}
and
\begin{equation}
\langle \widetilde{\phi}_{Im'}^{\mathbf{k}} |
\rho_{I,\nu} | u_{n\mathbf{k}} \rangle
=
- i \,
\langle
\partial_{k_\nu} \widetilde{\phi}_{Im'}^{\mathbf{k}}
| u_{n\mathbf{k}} \rangle .
\label{eq:rho_phi_right}
\end{equation}

Substituting into Eq.~\eqref{eq:hub_components},
\begin{align}
\Delta M_{\mathrm{Hub},\gamma}
&=
-\frac{\alpha}{2}
\sum_{n\mathbf{k}}^{\mathrm{occ}} f_{n\mathbf{k}}
\sum_I \sum_{mm'} \sum_{\mu\nu}
\epsilon_{\gamma\mu\nu}
\frac{1}{i}
\left(
i \langle u_{n\mathbf{k}} |
\partial_{k_\mu} \widetilde{\phi}_{Im}^{\mathbf{k}} \rangle
\right)
V_{mm'}^I
\left(
- i \langle
\partial_{k_\nu} \widetilde{\phi}_{Im'}^{\mathbf{k}}
| u_{n\mathbf{k}} \rangle
\right)
\nonumber\\
&=
\frac{i \alpha}{2}
\sum_{n\mathbf{k}}^{\mathrm{occ}} f_{n\mathbf{k}}
\sum_I \sum_{mm'} \sum_{\mu\nu}
\epsilon_{\gamma\mu\nu}
\langle u_{n\mathbf{k}} |
\partial_{k_\mu} \widetilde{\phi}_{Im}^{\mathbf{k}} \rangle
V_{mm'}^I
\langle
\partial_{k_\nu} \widetilde{\phi}_{Im'}^{\mathbf{k}}
| u_{n\mathbf{k}} \rangle .
\label{eq:hub_before_im}
\end{align}

Again, since the contraction with $\epsilon_{\gamma\mu\nu}$ selects the
antisymmetric part in the Cartesian indices, the result can be written
in manifestly real form as
\begin{equation}
\Delta M_{\mathrm{Hub},\gamma}
=
-\frac{\alpha}{2}
\sum_{n\mathbf{k}}^{\mathrm{occ}} f_{n\mathbf{k}}
\sum_I \sum_{mm'} \sum_{\mu\nu}
\epsilon_{\gamma\mu\nu}
\,\mathrm{Im}
\left[
\langle u_{n\mathbf{k}} |
\partial_{k_\mu} \widetilde{\phi}_{Im}^{\mathbf{k}} \rangle
V_{mm'}^I
\langle
\partial_{k_\nu} \widetilde{\phi}_{Im'}^{\mathbf{k}}
| u_{n\mathbf{k}} \rangle
\right] .
\label{eq:deltam_hub_final}
\end{equation}

In the present norm-conserving case, with \(S=\mathbb{I}\), the matrix \(V_{mm'}^I\) appearing in the separable Hubbard operator is identified with the on-site Hubbard potential matrix
\[
V_{mm'}^I \equiv U_{mm'}^{I,\sigma}
=
U^I \left[\frac{\delta_{mm'}}{2} - n_{mm'}^{I,\sigma}\right],
\]
with the spin index \(\sigma\) understood implicitly whenever not written explicitly.
\clearpage
The quantity
\[
\left\langle \partial_{k_\nu} \widetilde{\phi}_{Im'}^{\mathbf{k}} \middle| u_{n\mathbf{k}} \right\rangle
\]
is evaluated numerically by the routine \texttt{compute\_projU}, reported below, which is the direct analogue of \texttt{compute\_dbecp} for the Hubbard projectors (\texttt{wfcU} in QE). In the present implementation, the explicit stripping of the Bloch phase has not yet been introduced; however, tests indicate that this has only a minor impact on the numerical result.

\begin{verbatim}
  !------------------------------------------------------------------
  ! derivative of the wfcU (projU)
  !------------------------------------------------------------------
  SUBROUTINE compute_projU
    USE cell_base, ONLY : tpiba
    USE uspp_init, ONLY : init_us_2
    USE uspp,  ONLY : vkb
    USE ldaU, ONLY : wfcU
    implicit none
    integer :: ipol, sig
    integer :: iu
    real(dp) :: kq(3), xk_save(3)

    call deallocate_bec_type(becp)
    xk_save(:) = xk(:,ik)

    if (nkb == 0) return
    wfcU_save = wfcU

    do ipol = 1, 3
      projU(:,:,ipol) = (0.d0,0.d0)
      do sig = -1, 1, 2
        xk(:,ik) = xk_save(:)
        xk(ipol,ik) = xk(ipol,ik) + sig * delta_k
        call orthoUwfc_k( ik, .FALSE. )
        call strip_kphase_wfcU( ik )
        call calbec( npw, wfcU, evc, auxU, nbnd )
        projU(:,:,ipol) = projU(:,:,ipol) + &
                          0.5d0*sig/(delta_k*tpiba) * auxU(:,:)
      enddo
    enddo

    xk(:,ik) = xk_save(:)
    wfcU = wfcU_save
    CALL allocate_bec_type(nkb, nbnd, becp)
  END SUBROUTINE compute_projU
\end{verbatim}
The following routine removes the $k$-dependent Bloch phase from the Hubbard projectors before evaluating the non-local Hubbard contribution to the orbital magnetization.
\begin{verbatim}

SUBROUTINE strip_kphase_wfcU(ik)
  USE kinds,     ONLY : DP
  USE constants, ONLY : tpi          ! tpi = 2*pi in QE
  USE klist,     ONLY : xk
  USE ions_base, ONLY : nat, tau, ityp
  USE ldaU,      ONLY : wfcU, nwfcU, Hubbard_l  ! Hubbard_l(ntyp)
  IMPLICIT NONE

  INTEGER, INTENT(IN) :: ik
  INTEGER :: na, nt, l, m, iu, nproj
  REAL(DP) :: arg
  COMPLEX(DP) :: phase

  iu = 0

  DO na = 1, nat
     nt = ityp(na)
     l  = Hubbard_l(nt)

     ! Convenzione tipica: l < 0 => no Hubbard on this type
     IF (l >= 0) THEN
        nproj = 2*l + 1   ! collinear / NC

        arg = tpi * ( xk(1,ik)*tau(1,na) + xk(2,ik)*tau(2,na) + xk(3,ik)*tau(3,na) )
        phase = CMPLX( COS(arg), SIN(arg), KIND=DP )  ! exp(+i arg)

        DO m = 1, nproj
           iu = iu + 1
           wfcU(:,iu) = phase * wfcU(:,iu)
        END DO
     END IF
  END DO

  IF (iu /= nwfcU) THEN
     CALL errore('strip_kphase_wfcU', 'unexpected wfcU ordering: iu != nwfcU', 1)
  END IF
END SUBROUTINE strip_kphase_wfcU


\end{verbatim}
The final Hubbard contribution in Eq.~\eqref{eq:deltam_hub_final} is evaluated by the routine reported below:
\begin{verbatim}
  !------------------------------------------------------------------
  ! GIPAW correction (delta_M_Hub)
  !------------------------------------------------------------------
  SUBROUTINE calc_delta_M_hub
    USE ions_base,  ONLY : nat, ntyp => nsp, ityp
    USE ldaU,       ONLY : Hubbard_l, is_hubbard, offsetU, nwfcU
    USE scf,        ONLY : v
    implicit none
    complex(dp) :: tmp, hub_product
    integer :: ibnd, na, nt, m1, m2, ldim, off

    if (nwfcU == 0) return

    tmp = (0.d0, 0.d0)

    do nt = 1, ntyp
      if (.not. is_hubbard(nt)) cycle
      ldim = 2*Hubbard_l(nt) + 1

      do na = 1, nat
        if (ityp(na) /= nt) cycle
        off = offsetU(na)

        do m2 = 1, ldim
          do m1 = 1, ldim
            do ibnd = 1, nbnd
              hub_product = conjg(projU(off+m1, ibnd, ii)) * &
                                 projU(off+m2, ibnd, jj)
              tmp = tmp + wg(ibnd,ik) * &
                          v%ns(m1, m2, current_spin, na) * &
                          hub_product
            enddo
          enddo
        enddo
      enddo
    enddo

    delta_M_hub(kk) = delta_M_hub(kk) - 2.d0*imag(tmp)
  END SUBROUTINE calc_delta_M_hub
\end{verbatim}

For fixed Cartesian indices \(\mu \equiv ii\) and \(\nu \equiv jj\), this routine evaluates the quantity
\[
\sum_{n\mathbf{k}}^{\mathrm{occ}} w_{n\mathbf{k}}
\sum_I \sum_{mm'}
U_{mm'}^{I,\sigma}
\left[
\left\langle \partial_{k_\mu} \phi_{Im}^{\mathbf{k}} \middle| u_{n\mathbf{k}} \right\rangle
\right]^*
\left\langle \partial_{k_\nu} \phi_{Im'}^{\mathbf{k}} \middle| u_{n\mathbf{k}} \right\rangle ,
\]
where the Hubbard potential matrix \(U_{mm'}^{I,\sigma}\) is stored as
\texttt{v\%ns(m1,m2,current\_spin,na)}, while the derivative matrix elements are provided by \texttt{projU}. The corresponding contribution to the \(\kk\)-th component of the magnetization is then obtained as
\[
\Delta M_{\mathrm{Hub},\kk} \propto -2\,\mathrm{Im}\!\left[\cdots\right].
\]



\subsection{Numerical test with LiCoO$_2$.}

As a numerical benchmark, we consider the NMR shielding tensor of
LiCoO$_2$ at the DFT+$U$ level, comparing QE-CONVERSE with the
linear-response GIPAW implementation in CASTEP. Norm-conserving
pseudopotentials and a plane-wave cutoff of 60 Ry were used, with a
Hubbard correction of $U=8$ eV applied to the Co-$3d$ manifold. The
comparison with CASTEP is particularly relevant, since its GIPAW
implementation explicitly includes the non-local Hubbard contribution in
the magnetic-response operator.

Table~\ref{tab:licoo2_dftu_shielding} reports the isotropic shielding
values for the three equivalent Li, Co, and O atoms in the unit cell. For
each atomic species, both methods give internally consistent values across
the three symmetry-equivalent sites. The agreement is satisfactory for Li
and Co, whereas a sizeable discrepancy is observed for the oxygen
shielding. This is consistent with the fact that, although the Hubbard
correction is applied only to the Co-$3d$ manifold, the oxygen shielding is
sensitive to the Co-$3d$/O-$2p$ hybridization and therefore to the detailed
definition of the Hubbard subspace. A possible contribution to the
remaining discrepancy is the different construction of the Hubbard
projectors and of the associated non-local magnetic-response term in the
two implementations.
\begin{table}[h]
\centering
\caption{$^{7}$Li, $^{59}$Co, and $^{17}$O isotropic NMR shielding (ppm) in LiCoO$_2$ at the DFT+$U$ level ($U=8$ eV on Co-$3d$), computed with QE-CONVERSE and CASTEP. Three symmetry-equivalent atoms of each species are reported.}
\label{tab:licoo2_dftu_shielding}
\begin{tabular}{llcc}
\hline
Atom & Site & QE-CONVERSE (ppm) & CASTEP (ppm) \\
\hline
Li & 1 & 92.51 & 87.06 \\
Li & 2 & 92.33 & 87.06 \\
Li & 3 & 92.32 & 87.06 \\
\hline
Co & 1 & -4679.62 & -4807.13 \\
Co & 2 & -4679.64 & -4807.12 \\
Co & 3 & -4679.64 & -4807.11 \\
\hline
O & 1 & 343.12 & 230.66 \\
O & 2 & 343.11 & 230.68 \\
O & 3 & 343.03 & 230.67 \\
\hline
\end{tabular}
\end{table}

\begin{table}[h]
\centering
\caption{Full NMR shielding tensor components and isotropic shieldings (ppm) for all atoms in LiCoO$_2$ at the DFT+$U$+$V$ level, computed with QE-CONVERSE. $\sigma_\text{iso} = (\sigma_{xx}+\sigma_{yy}+\sigma_{zz})/3$.}
\label{tab:licoo2_dftu_shielding_full}
\begin{tabular}{lrrrr}
\hline
Atom & $\sigma_{xx}$ & $\sigma_{yy}$ & $\sigma_{zz}$ & $\sigma_\text{iso}$ \\
\hline
Co 1 & $-4982.31$ & $-4982.89$ & $-5367.68$ & $-5110.96$ \\
Co 2 & $-4982.80$ & $-4983.56$ & $-5367.89$ & $-5111.41$ \\
Co 3 & $-4982.26$ & $-4982.93$ & $-5367.91$ & $-5111.03$ \\
\hline
O 4  & $66.37$    & $66.26$    & $65.90$    & $66.18$    \\
O 5  & $66.69$    & $66.42$    & $66.15$    & $66.42$    \\
O 6  & $66.37$    & $66.30$    & $66.20$    & $66.29$    \\
O 7  & $66.33$    & $66.32$    & $66.19$    & $66.28$    \\
O 8  & $66.65$    & $66.49$    & $66.06$    & $66.40$    \\
O 9  & $66.39$    & $66.14$    & $65.94$    & $66.16$    \\
\hline
Li 10 & $88.41$   & $88.27$    & $82.43$    & $86.37$    \\
Li 11 & $88.27$   & $88.23$    & $82.61$    & $86.37$    \\
Li 12 & $88.24$   & $88.24$    & $82.65$    & $86.38$    \\
\hline
\end{tabular}
\end{table}

\newpage
\clearpage

\section{DFT+$U$+$V$}

In DFT+$U$+$V$, a corrective Hubbard energy $E_{U+V}$ is added to the standard DFT energy $E_{\mathrm{DFT}}$:
%
\begin{equation}
    E_{\mathrm{DFT}+U+V} = E_{\mathrm{DFT}} + E_{U+V} .
    \label{eq:Edft_plus_uv}
\end{equation}
%
Unlike DFT+$U$, which only accounts for on-site interactions scaled by $U$, DFT+$U$+$V$ also includes inter-site interactions between an atom and its neighbors, scaled by $V$. In the simplified rotationally-invariant formulation, the extended Hubbard correction reads:
%
\begin{eqnarray}
    E_{U+V} & = & \frac{1}{2} \sum_I \sum_{\sigma m m'}
    U^I \left( \delta_{m m'} - n^{II \sigma}_{m m'} \right) n^{II \sigma}_{m' m} \nonumber \\
    %
    & & - \frac{1}{2} \sum_{I} \sum_{J (J \ne I)}^{*} \sum_{\sigma m m'} V^{I J}
    n^{I J \sigma}_{m m'} n^{J I \sigma}_{m' m} \,,
    \label{eq:Edftu}
\end{eqnarray}
%
where $I$ and $J$ index atomic sites, $m$ and $m'$ denote magnetic quantum numbers, and $U^I$ and $V^{IJ}$ are the \textit{effective} on-site and inter-site Hubbard parameters, respectively (here and everywhere in the following, $U^I$ denotes the difference between the screened Coulomb interaction parameter and the Hund's exchange interaction parameter). The asterisk denotes that, for each atom $I$, the sum over $J$ includes all neighbors up to a defined cutoff distance. The generalized occupation matrices $n^{IJ \sigma}_{m m'}$ are obtained by projecting KS states $\psi^\sigma_{v,\mathbf{k}}$ onto localized orbitals $\phi^I_m(\mathbf{r})$:
%
\begin{equation}
    n^{I J \sigma}_{m m'} = \sum_{v,\mathbf{k}} f^\sigma_{v,\mathbf{k}}
    \braket{\psi^\sigma_{v,\mathbf{k}}}{\phi^{J}_{m'}} \braket{\phi^{I}_{m}}{\psi^\sigma_{v,\mathbf{k}}} \,,
    \label{eq:occ_matrix_0}
\end{equation}
%
where $v$ and $\sigma$ label bands and spin, $\mathbf{k}$ runs over the first Brillouin zone (BZ), and $f^\sigma_{v,\mathbf{k}}$ are the KS occupations. In Eq.~\eqref{eq:occ_matrix_0}, $m \in I$ and $m' \in J$, meaning that $m$ runs over the magnetic quantum numbers associated with a fixed principal quantum number $n$ and orbital angular momentum $l$ of atom $I$, and analogously $m'$ for atom $J$. The corresponding generalized KS equation can be written as,
%
\begin{equation}
     \left(\hat{\mathcal H}_\mathrm{KS}^\sigma +
     \hat{\mathcal V}_{U+V}^\sigma \right)|\psi_{v,{\bf k}}^\sigma\rangle
     = \varepsilon_{v,{\bf k}}^\sigma |\psi_{v,{\bf k}}^\sigma\rangle,
\end{equation}
%
where $\hat{\mathcal H}_\mathrm{KS}^\sigma$ is the KS Hamiltonian of the standard DFT, $\varepsilon_{v,{\bf k}}^\sigma$ is the $v$-th band energy at momentum ${\bf k}$ and spin $\sigma$, and $\hat{\mathcal V}_{U+V}^\sigma$ is the Hubbard potential that is derived from the Hubbard energy~\eqref{eq:Edftu} and which reads:
%
\begin{eqnarray} 
\hat{\mathcal V}_{U+V}^\sigma & = &
    \sum_{I}\sum_{mm'}\frac{U^I}{2}
    \left(
    \delta_{mm'}-2n^{II\sigma}_{mm'}
    \right) 
    |\phi^I_m\rangle
    \langle\phi^I_{m'}| \nonumber \\
    % 
    & & - \sum_{I} \sum_{J (J \ne I)}^{*}\sum_{mm'} V^{IJ} n^{IJ\sigma}_{mm'}
    |\phi^I_m\rangle
    \langle \phi^J_{m'}|.
    \label{eq:Vdftuv}
\end{eqnarray}
%
In Eqs.~\eqref{eq:Edftu} and \eqref{eq:Vdftuv}, the effective on-site $U^I$ and inter-site $V^{IJ}$ terms act in opposition when both parameters are positive: the on-site term favors electron localization on atomic sites and suppresses hybridization, whereas the inter-site term enhances hybridization with neighboring atoms. Accurate first-principles evaluation of $U^I$ and $V^{IJ}$ is therefore crucial to correctly describe the balance between localization and hybridization. 

It is important to note that the computed values of $U^I$ and $V^{IJ}$ depend on the choice of Hubbard projector functions $\phi^I_{m}(\mathbf{r})$ used to define the occupation matrix in Eq.~\eqref{eq:occ_matrix_0} and the Hubbard potential in Eq.~\eqref{eq:Vdftuv}. 

\subsection{Contribution to the orbital magnetization}

We now derive the contribution of the full DFT$+U+V$ potential to the orbital
magnetization, generalizing the DFT$+U$ result of
Section~\ref{sec:hubmagnetization}.

\paragraph{Notation.}
It is convenient to introduce a unified potential matrix that covers both the
on-site ($J=I$) and inter-site ($J\ne I$) terms of
Eq.~\eqref{eq:Vdftuv}:
\begin{equation}
  W_{mm'}^{IJ,\sigma}
  =
  \begin{cases}
    \dfrac{U^I}{2}\!\left(\delta_{mm'} - 2n_{mm'}^{II\sigma}\right)
    & J = I, \\[6pt]
    -V^{IJ}\,n_{mm'}^{IJ\sigma}
    & J \ne I.
  \end{cases}
  \label{eq:Wmat}
\end{equation}
In QE the full matrix $W_{mm'}^{IJ,\sigma}$ is stored as
\texttt{v\_nsg(m,m',viz,na,$\sigma$)}, where \texttt{na}$=I$ and
\texttt{viz} is the neighbor index such that
\texttt{neighood(na)\%neigh(viz)}$=J$.  The on-site entry corresponds to
the self-neighbor (\texttt{neigh(viz)}$=I$), so the sum over all viz
automatically includes both parts.

Using $W_{mm'}^{IJ}$, the DFT$+U+V$ potential reads
\begin{equation}
  \hat{\mathcal{V}}_{U+V}^{\sigma}
  =
  \sum_{I}\sum_{J}^{*}\sum_{mm'}
  W_{mm'}^{IJ,\sigma}\,
  |\phi^{I}_{m}\rangle\langle\phi^{J}_{m'}|,
  \label{eq:Vuv_compact}
\end{equation}
where the asterisk on $\sum_J$ again denotes the sum over the self-site
($J=I$) and all neighbors within the cutoff.

\paragraph{On-site part ($J=I$).}
The on-site contribution has the same operator structure as the
DFT$+U$ potential: both bra and ket belong to atom $I$.  The
derivation of Section~\ref{sec:hubmagnetization} therefore applies
unchanged, with $V_{mm'}^I \to W_{mm'}^{II,\sigma}$.  It yields
Eq.~\eqref{eq:deltam_hub_final} with this replacement.

\paragraph{Inter-site part ($J\ne I$).}
For the inter-site terms, the ket $|\phi^I_m\rangle$ belongs to atom
$I$ but the bra $\langle\phi^J_{m'}|$ belongs to atom $J\ne I$.
Starting from Eq.~\eqref{eq:generic_form}, the inter-site contribution
to the component $\gamma$ of the orbital magnetization reads
\begin{equation}
  \Delta M_{V,\gamma}
  =
  -\frac{\alpha}{2}
  \sum_{n\mathbf{k}}^{\mathrm{occ}}f_{n\mathbf{k}}
  \sum_{I}\sum_{J\ne I}^{*}\sum_{mm'}\sum_{\mu\nu}
  \epsilon_{\gamma\mu\nu}\,\frac{1}{i}
  \langle u_{n\mathbf{k}}|\rho_{I,\mu}|\widetilde\phi^{\mathbf{k},I}_{m}\rangle
  \,W_{mm'}^{IJ,\sigma}\,
  \langle\widetilde\phi^{\mathbf{k},J}_{m'}|\rho_{I,\nu}|u_{n\mathbf{k}}\rangle.
  \label{eq:deltaM_V_components}
\end{equation}
The left matrix element is handled exactly as in the DFT$+U$ case via
the no-phase identity~\eqref{eq:rho_phi_left}:
\begin{equation}
  \langle u_{n\mathbf{k}}|\rho_{I,\mu}|\widetilde\phi^{\mathbf{k},I}_{m}\rangle
  =
  i\,\langle u_{n\mathbf{k}}|\partial_{k_\mu}\widetilde\phi^{\mathbf{k},I}_{m}\rangle.
  \label{eq:left_I}
\end{equation}

For the right matrix element we decompose the relative coordinate with
respect to $I$ into the relative coordinate with respect to $J$ plus
the inter-site displacement $\mathbf{d}^{IJ}=\mathbf{R}_J-\mathbf{R}_I$:
\begin{equation}
  \rho_{I,\nu} = \rho_{J,\nu} + d^{IJ}_\nu,
  \qquad
  \mathbf{d}^{IJ} = \mathbf{R}_J - \mathbf{R}_I.
  \label{eq:rho_decomp}
\end{equation}
Applying the no-phase identity for atom $J$,
Eq.~\eqref{eq:rho_phi_right}, to the first term, and collecting:
\begin{equation}
  \langle\widetilde\phi^{\mathbf{k},J}_{m'}|\rho_{I,\nu}|u_{n\mathbf{k}}\rangle
  =
  -i\,\langle\partial_{k_\nu}\widetilde\phi^{\mathbf{k},J}_{m'}|u_{n\mathbf{k}}\rangle
  +
  d^{IJ}_\nu\,
  \langle\widetilde\phi^{\mathbf{k},J}_{m'}|u_{n\mathbf{k}}\rangle.
  \label{eq:right_IJ}
\end{equation}
The second term is new: it involves the regular (non-derivative) Hubbard
overlap
\begin{equation}
  s^{J}_{m',n}(\mathbf{k})
  \;\equiv\;
  \langle\widetilde\phi^{\mathbf{k},J}_{m'}|u_{n\mathbf{k}}\rangle,
  \label{eq:hubbecp}
\end{equation}
which is the no-phase analogue of the standard $\langle\beta|u\rangle$
projector overlap.

Substituting Eqs.~\eqref{eq:left_I}--\eqref{eq:right_IJ} into
Eq.~\eqref{eq:deltaM_V_components}:
\begin{align}
  \Delta M_{V,\gamma}
  &=
  -\frac{\alpha}{2}
  \sum_{n\mathbf{k}}^{\mathrm{occ}}f_{n\mathbf{k}}
  \sum_{I,J\ne I}^{*}\sum_{mm'}\sum_{\mu\nu}
  \epsilon_{\gamma\mu\nu}\,
  \langle u_{n\mathbf{k}}|\partial_{k_\mu}\widetilde\phi^{\mathbf{k},I}_{m}\rangle
  W_{mm'}^{IJ,\sigma}
  \!\left[
    -i\,\langle\partial_{k_\nu}\widetilde\phi^{\mathbf{k},J}_{m'}|u_{n\mathbf{k}}\rangle
    +
    d^{IJ}_\nu\,s^{J}_{m',n}
  \right].
  \label{eq:deltaM_V_split}
\end{align}
This splits into two structurally distinct contributions.

\paragraph{Part~A: derivative--derivative term.}
The first bracket term gives
\begin{equation}
  \Delta M_{V,\gamma}^{A}
  =
  \frac{i\alpha}{2}
  \sum_{n\mathbf{k}}^{\mathrm{occ}}f_{n\mathbf{k}}
  \sum_{I,J\ne I}^{*}\sum_{mm'}\sum_{\mu\nu}
  \epsilon_{\gamma\mu\nu}\,
  \langle u_{n\mathbf{k}}|\partial_{k_\mu}\widetilde\phi^{\mathbf{k},I}_{m}\rangle
  W_{mm'}^{IJ,\sigma}\,
  \langle\partial_{k_\nu}\widetilde\phi^{\mathbf{k},J}_{m'}|u_{n\mathbf{k}}\rangle,
  \label{eq:deltaM_VA}
\end{equation}
which has exactly the same structure as the DFT$+U$ result
Eq.~\eqref{eq:hub_before_im}, now with projector derivatives from two
\emph{different} atomic centers.  Passing to the imaginary part as in
Section~\ref{sec:hubmagnetization}:
\begin{equation}
  \Delta M_{V,\gamma}^{A}
  =
  -\frac{\alpha}{2}
  \sum_{n\mathbf{k}}^{\mathrm{occ}}f_{n\mathbf{k}}
  \sum_{I,J\ne I}^{*}\sum_{mm'}\sum_{\mu\nu}
  \epsilon_{\gamma\mu\nu}\,
  \mathrm{Im}\!\left[
    \langle u_{n\mathbf{k}}|\partial_{k_\mu}\widetilde\phi^{\mathbf{k},I}_{m}\rangle
    W_{mm'}^{IJ,\sigma}\,
    \langle\partial_{k_\nu}\widetilde\phi^{\mathbf{k},J}_{m'}|u_{n\mathbf{k}}\rangle
  \right].
  \label{eq:deltaM_VA_im}
\end{equation}

\paragraph{Part~B: derivative--overlap term.}
The second bracket term gives
\begin{equation}
  \Delta M_{V,\gamma}^{B}
  =
  -\frac{\alpha}{2}
  \sum_{n\mathbf{k}}^{\mathrm{occ}}f_{n\mathbf{k}}
  \sum_{I,J\ne I}^{*}\sum_{mm'}\sum_{\mu\nu}
  \epsilon_{\gamma\mu\nu}\,d^{IJ}_\nu\,
  \langle u_{n\mathbf{k}}|\partial_{k_\mu}\widetilde\phi^{\mathbf{k},I}_{m}\rangle
  W_{mm'}^{IJ,\sigma}\,
  s^{J}_{m',n}.
  \label{eq:deltaM_VB}
\end{equation}
This term is absent in the pure DFT$+U$ case (where $J=I$ so
$\mathbf{d}^{II}=\mathbf{0}$) and arises solely from the inter-site
nature of the projector pair. Since $d^{IJ}_\nu$ is real while
$W_{mm'}^{IJ,\sigma}$ and $s^J_{m',n}$ are in general complex, the
expression under the sum is complex; the physical result is its real
part (the imaginary parts cancel in the full sum over ordered pairs
$(I,J)$ and $(J,I)$, as follows from $W^{JI}_{m'm} = (W^{IJ}_{mm'})^*$
and time-reversal):
\begin{equation}
  \Delta M_{V,\gamma}^{B}
  =
  -\frac{\alpha}{2}
  \sum_{n\mathbf{k}}^{\mathrm{occ}}f_{n\mathbf{k}}
  \sum_{I,J\ne I}^{*}\sum_{mm'}\sum_{\mu\nu}
  \epsilon_{\gamma\mu\nu}\,d^{IJ}_\nu\,
  \mathrm{Re}\!\left[
    \langle u_{n\mathbf{k}}|\partial_{k_\mu}\widetilde\phi^{\mathbf{k},I}_{m}\rangle
    W_{mm'}^{IJ,\sigma}\,
    s^{J}_{m',n}
  \right].
  \label{eq:deltaM_VB_re}
\end{equation}

\paragraph{Full DFT$+U+V$ magnetization formula.}
Combining the on-site part (Eq.~\eqref{eq:deltam_hub_final} with
$V^I\to W^{II}$) with Parts~A and B, the total contribution reads
\begin{align}
  \Delta M_{U+V,\gamma}
  &=
  -\frac{\alpha}{2}
  \sum_{n\mathbf{k}}^{\mathrm{occ}}f_{n\mathbf{k}}
  \sum_{I}\sum_{J}^{*}\sum_{mm'}\sum_{\mu\nu}
  \epsilon_{\gamma\mu\nu}\,
  \mathrm{Im}\!\left[
    \langle u_{n\mathbf{k}}|\partial_{k_\mu}\widetilde\phi^{\mathbf{k},I}_{m}\rangle
    W_{mm'}^{IJ,\sigma}\,
    \langle\partial_{k_\nu}\widetilde\phi^{\mathbf{k},J}_{m'}|u_{n\mathbf{k}}\rangle
  \right]
  \nonumber\\
  &\quad
  -\frac{\alpha}{2}
  \sum_{n\mathbf{k}}^{\mathrm{occ}}f_{n\mathbf{k}}
  \sum_{I}\sum_{J\ne I}^{*}\sum_{mm'}\sum_{\mu\nu}
  \epsilon_{\gamma\mu\nu}\,d^{IJ}_\nu\,
  \mathrm{Re}\!\left[
    \langle u_{n\mathbf{k}}|\partial_{k_\mu}\widetilde\phi^{\mathbf{k},I}_{m}\rangle
    W_{mm'}^{IJ,\sigma}\,
    s^{J}_{m',n}
  \right].
  \label{eq:deltam_uv_final}
\end{align}
The first line (the Im term) extends over \emph{all} neighbors including
$J=I$; for $J=I$ it reduces to Eq.~\eqref{eq:deltam_hub_final} since
$W^{II}=V^{II}$ (the DFT$+U$ potential matrix).  The second line
(the Re term) is strictly inter-site and vanishes when $V=0$.

\paragraph{Implementation notes.}
\begin{itemize}
  \item The derivative overlaps
    $\langle u_{n\mathbf{k}}|\partial_{k_\mu}\widetilde\phi^{\mathbf{k},I}_{m}\rangle$
    and
    $\langle\partial_{k_\nu}\widetilde\phi^{\mathbf{k},J}_{m'}|u_{n\mathbf{k}}\rangle$
    are already computed by \texttt{compute\_dhubbecp} for every Hubbard
    atom.

  \item The regular overlap $s^J_{m',n}=\langle\widetilde\phi^{\mathbf{k},J}_{m'}|u_{n\mathbf{k}}\rangle$
    must be computed separately, e.g.\ by calling \texttt{orthoUwfc\_k}
    at the unshifted $\mathbf{k}$, stripping the phase via
    \texttt{strip\_kphase\_wfcU}, and projecting with \texttt{calbec}.

  \item The potential matrix $W_{mm'}^{IJ,\sigma}$ is stored as
    \texttt{v\_nsg(m,m',viz,na,$\sigma$)}, where \texttt{na}$=I$ and
    \texttt{neighood(na)\%neigh(viz)}$=J$.

  \item The inter-site displacement $\mathbf{d}^{IJ}=\mathbf{R}_J-\mathbf{R}_I$
    is available from \texttt{tau(:,na2)~-~tau(:,na1)} in crystal
    coordinates (converted to Cartesian with \texttt{alat}).
\end{itemize}


\appendix
\section*{Appendix}
\addcontentsline{toc}{section}{Appendix}

\section{Phase convention for the $k$-derivative of Kleinman--Bylander projectors}
\label{app:kb_no_phase}

In this Appendix we discuss the phase convention used to compute the
$k$-derivative of Kleinman--Bylander (KB) projectors. The purpose is to
show that the derivative entering the GIPAW/converse expression must be
taken with respect to the local coordinate measured from the atomic
center. This requires removing the global phase associated with the
position of the atomic center.

Let $\mathbf r$ denote the electronic coordinate and let $\mathbf R_I$
be the position of atom $I$. We define the local coordinate with respect
to atom $I$ as
\begin{equation}
\boldsymbol{\rho}_I = \mathbf r - \mathbf R_I .
\label{eq:rho_def_app}
\end{equation}
The localized projector centered on atom $I$ is denoted by
\begin{equation}
\beta_i^I(\mathbf r) \equiv \beta_i(\mathbf r-\mathbf R_I)
= \beta_i(\boldsymbol{\rho}_I) .
\label{eq:beta_centered_def}
\end{equation}

The GIPAW contribution involving a non-local operator centered on atom
$I$ contains the relative position operator
\begin{equation}
\mathbf r - \mathbf R_I = \boldsymbol{\rho}_I ,
\end{equation}
rather than the absolute position $\mathbf r$. Therefore, the
$k$-derivative of the projector must reproduce the action of
$\boldsymbol{\rho}_I$.

\subsection{Full-phase projector}

Consider first the projector with the full Bloch phase. In real space,
this object may be written as
\begin{equation}
\beta_{i,\mathrm{full}}^{\mathbf k,I}(\mathbf r)
=
e^{-i\mathbf k\cdot\mathbf r}\,
\beta_i^I(\mathbf r) .
\label{eq:beta_full_real}
\end{equation}
Using Eq.~\eqref{eq:beta_centered_def}, this can also be written as
\begin{equation}
\beta_{i,\mathrm{full}}^{\mathbf k,I}(\mathbf r)
=
e^{-i\mathbf k\cdot\mathbf r}\,
\beta_i(\mathbf r-\mathbf R_I) .
\end{equation}
Taking the derivative with respect to $k_\mu$ gives
\begin{equation}
\partial_{k_\mu}
\beta_{i,\mathrm{full}}^{\mathbf k,I}(\mathbf r)
=
-i r_\mu\,
\beta_{i,\mathrm{full}}^{\mathbf k,I}(\mathbf r) .
\label{eq:full_der_real}
\end{equation}
Thus, the full-phase derivative generates the absolute coordinate
$r_\mu$, not the relative coordinate $r_\mu-R_{I,\mu}$.

The same result can be seen in reciprocal space. The plane-wave
coefficients of the full-phase projector have the structure
\begin{equation}
\beta_{i,\mathrm{full}}^{\mathbf k,I}(\mathbf G)
=
e^{-i(\mathbf k+\mathbf G)\cdot\mathbf R_I}
\beta_i(\mathbf k+\mathbf G),
\label{eq:beta_full_g}
\end{equation}
where $\beta_i(\mathbf k+\mathbf G)$ is the Fourier transform of the
projector centered at the origin. Separating the phase factor gives
\begin{equation}
\beta_{i,\mathrm{full}}^{\mathbf k,I}(\mathbf G)
=
e^{-i\mathbf k\cdot\mathbf R_I}
e^{-i\mathbf G\cdot\mathbf R_I}
\beta_i(\mathbf k+\mathbf G).
\label{eq:beta_full_g_sep}
\end{equation}
Taking the derivative with respect to $k_\mu$ yields
\begin{align}
\partial_{k_\mu}
\beta_{i,\mathrm{full}}^{\mathbf k,I}(\mathbf G)
&=
\partial_{k_\mu}
\left[
e^{-i\mathbf k\cdot\mathbf R_I}
e^{-i\mathbf G\cdot\mathbf R_I}
\beta_i(\mathbf k+\mathbf G)
\right]
\nonumber\\
&=
e^{-i\mathbf k\cdot\mathbf R_I}
e^{-i\mathbf G\cdot\mathbf R_I}
\partial_{k_\mu}\beta_i(\mathbf k+\mathbf G)
\nonumber\\
&\quad
-iR_{I,\mu}
e^{-i\mathbf k\cdot\mathbf R_I}
e^{-i\mathbf G\cdot\mathbf R_I}
\beta_i(\mathbf k+\mathbf G).
\label{eq:full_der_g}
\end{align}
The second term in Eq.~\eqref{eq:full_der_g} is explicitly proportional
to the atomic position $R_{I,\mu}$. This contribution originates from
the derivative of the global phase
\begin{equation}
e^{-i\mathbf k\cdot\mathbf R_I}.
\end{equation}
It is not associated with the intrinsic spatial dependence of the
localized projector, but only with the absolute position of its center.

This term is undesirable for the GIPAW/converse expression, because the
operator entering the commutator is $\mathbf r-\mathbf R_I$, not
$\mathbf r$. Equivalently, the derivative must be insensitive to an
arbitrary global phase attached to the atomic center. Keeping the
full phase would introduce an origin-dependent contribution proportional
to $\mathbf R_I$.

\subsection{No-phase projector}

To remove this spurious contribution, one defines the no-phase projector
by dropping the factor $e^{-i\mathbf k\cdot\mathbf R_I}$ from
Eq.~\eqref{eq:beta_full_g_sep}. In reciprocal space, the no-phase
projector is therefore
\begin{equation}
\widetilde{\beta}_{i}^{\mathbf k,I}(\mathbf G)
=
e^{-i\mathbf G\cdot\mathbf R_I}
\beta_i(\mathbf k+\mathbf G).
\label{eq:beta_no_phase_g}
\end{equation}
This is the object constructed by the no-phase version of the KB
projector initialization. In the code, the factor
$e^{-i\mathbf G\cdot\mathbf R_I}$ is the structure factor associated
with the reciprocal lattice vector $\mathbf G$, whereas the phase
$e^{-i\mathbf k\cdot\mathbf R_I}$ is deliberately omitted.

Taking the derivative of Eq.~\eqref{eq:beta_no_phase_g} gives
\begin{equation}
\partial_{k_\mu}
\widetilde{\beta}_{i}^{\mathbf k,I}(\mathbf G)
=
e^{-i\mathbf G\cdot\mathbf R_I}
\partial_{k_\mu}\beta_i(\mathbf k+\mathbf G).
\label{eq:no_phase_der_g}
\end{equation}
Unlike Eq.~\eqref{eq:full_der_g}, there is no term proportional to
$R_{I,\mu}$. Thus, the derivative acts only on the intrinsic dependence
of the projector on $\mathbf k+\mathbf G$.

\subsection{Real-space reconstruction of the no-phase projector}

We now show explicitly that the no-phase reciprocal-space definition
corresponds to a real-space projector whose $k$-derivative generates
the relative coordinate $\mathbf r-\mathbf R_I$.

The inverse plane-wave expansion of Eq.~\eqref{eq:beta_no_phase_g} is
\begin{equation}
\widetilde{\beta}_{i}^{\mathbf k,I}(\mathbf r)
=
\frac{1}{\sqrt{\Omega}}
\sum_{\mathbf G}
\widetilde{\beta}_{i}^{\mathbf k,I}(\mathbf G)
e^{i\mathbf G\cdot\mathbf r},
\label{eq:inverse_pw_no_phase}
\end{equation}
where $\Omega$ is the cell volume. Substituting
Eq.~\eqref{eq:beta_no_phase_g}, one obtains
\begin{align}
\widetilde{\beta}_{i}^{\mathbf k,I}(\mathbf r)
&=
\frac{1}{\sqrt{\Omega}}
\sum_{\mathbf G}
e^{-i\mathbf G\cdot\mathbf R_I}
\beta_i(\mathbf k+\mathbf G)
e^{i\mathbf G\cdot\mathbf r}
\nonumber\\
&=
\frac{1}{\sqrt{\Omega}}
\sum_{\mathbf G}
\beta_i(\mathbf k+\mathbf G)
e^{i\mathbf G\cdot(\mathbf r-\mathbf R_I)} .
\label{eq:no_phase_real_sum}
\end{align}
Using the local coordinate
$\boldsymbol{\rho}_I=\mathbf r-\mathbf R_I$, this becomes
\begin{equation}
\widetilde{\beta}_{i}^{\mathbf k,I}(\mathbf r)
=
\frac{1}{\sqrt{\Omega}}
\sum_{\mathbf G}
\beta_i(\mathbf k+\mathbf G)
e^{i\mathbf G\cdot\boldsymbol{\rho}_I}.
\label{eq:no_phase_real_rho}
\end{equation}

The quantity $\beta_i(\mathbf k+\mathbf G)$ is the Fourier coefficient
of the localized function multiplied by the Bloch phase relative to the
atomic center. Hence Eq.~\eqref{eq:no_phase_real_rho} is the
plane-wave representation of
\begin{equation}
\widetilde{\beta}_{i}^{\mathbf k,I}(\mathbf r)
=
e^{-i\mathbf k\cdot(\mathbf r-\mathbf R_I)}
\beta_i(\mathbf r-\mathbf R_I),
\label{eq:no_phase_real_final}
\end{equation}
up to the usual periodic repetition implied by the plane-wave
representation.

Taking the derivative of Eq.~\eqref{eq:no_phase_real_final} gives
\begin{align}
\partial_{k_\mu}
\widetilde{\beta}_{i}^{\mathbf k,I}(\mathbf r)
&=
\partial_{k_\mu}
\left[
e^{-i\mathbf k\cdot(\mathbf r-\mathbf R_I)}
\beta_i(\mathbf r-\mathbf R_I)
\right]
\nonumber\\
&=
-i(r_\mu-R_{I,\mu})
e^{-i\mathbf k\cdot(\mathbf r-\mathbf R_I)}
\beta_i(\mathbf r-\mathbf R_I)
\nonumber\\
&=
-i(r_\mu-R_{I,\mu})
\widetilde{\beta}_{i}^{\mathbf k,I}(\mathbf r)
\nonumber\\
&=
-i\rho_{I,\mu}
\widetilde{\beta}_{i}^{\mathbf k,I}(\mathbf r).
\label{eq:no_phase_der_real}
\end{align}
Therefore,
\begin{equation}
(r_\mu-R_{I,\mu})
\widetilde{\beta}_{i}^{\mathbf k,I}(\mathbf r)
=
i\,
\partial_{k_\mu}
\widetilde{\beta}_{i}^{\mathbf k,I}(\mathbf r).
\label{eq:rho_der_identity}
\end{equation}

Equation~\eqref{eq:rho_der_identity} is the identity required in the
GIPAW/converse derivation. It shows that the no-phase derivative
generates the relative position operator with respect to the atomic
center.

\subsection{No-phase convention for Hubbard projectors}
\label{app:hubbard_no_phase}

The same phase convention must be used for the Hubbard projectors. Let
$\phi_{Im}$ be a localized Hubbard orbital centered on atom $I$. In
reciprocal space, the full-phase Hubbard projector has the same
translation structure as a KB projector:
\begin{equation}
\phi^{\mathbf k,I}_{m,\mathrm{full}}(\mathbf G)
=
e^{-i(\mathbf k+\mathbf G)\cdot\mathbf R_I}
\phi_m(\mathbf k+\mathbf G).
\end{equation}
Separating the phase associated with the atomic center gives
\begin{equation}
\phi^{\mathbf k,I}_{m,\mathrm{full}}(\mathbf G)
=
e^{-i\mathbf k\cdot\mathbf R_I}
e^{-i\mathbf G\cdot\mathbf R_I}
\phi_m(\mathbf k+\mathbf G).
\end{equation}
The factor $e^{-i\mathbf G\cdot\mathbf R_I}$ must be retained, because
it places the localized orbital on atom $I$. By contrast, the factor
$e^{-i\mathbf k\cdot\mathbf R_I}$ is a global, $k$-dependent phase
associated with the position of the atomic center. If differentiated
with respect to $\mathbf k$, it generates an additional term
proportional to $\mathbf R_I$.

Therefore we define the no-phase Hubbard projector as
\begin{equation}
\widetilde{\phi}^{\mathbf k,I}_{m}(\mathbf G)
=
e^{+i\mathbf k\cdot\mathbf R_I}
\phi^{\mathbf k,I}_{m,\mathrm{full}}(\mathbf G),
\end{equation}
or equivalently
\begin{equation}
\widetilde{\phi}^{\mathbf k,I}_{m}(\mathbf G)
=
e^{-i\mathbf G\cdot\mathbf R_I}
\phi_m(\mathbf k+\mathbf G).
\end{equation}

With this definition, the inverse plane-wave expansion gives the
cell-periodic projector
\begin{equation}
\widetilde{\phi}^{\mathbf k,I}_{m}(\mathbf r)
=
e^{-i\mathbf k\cdot(\mathbf r-\mathbf R_I)}
\phi_m(\mathbf r-\mathbf R_I),
\end{equation}
up to the periodic repetition implied by the plane-wave representation.
Consequently,
\begin{equation}
\partial_{k_\mu}
\widetilde{\phi}^{\mathbf k,I}_{m}(\mathbf r)
=
-i(r_\mu-R_{I,\mu})
\widetilde{\phi}^{\mathbf k,I}_{m}(\mathbf r).
\end{equation}
Thus,
\begin{equation}
(r_\mu-R_{I,\mu})
\widetilde{\phi}^{\mathbf k,I}_{m}
=
i\,\partial_{k_\mu}
\widetilde{\phi}^{\mathbf k,I}_{m}.
\end{equation}

This is the same identity used for the no-phase KB projectors. Therefore,
when the Hubbard contribution is written in terms of
$\partial_{\mathbf k}\phi_{Im}^{\mathbf k}$, the derivative must be
taken on the no-phase Hubbard projectors. Otherwise the derivative would
also act on the global phase $e^{-i\mathbf k\cdot\mathbf R_I}$ and would
produce a spurious contribution proportional to $R_{I,\mu}$.


\end{document}



